\documentclass[12pt]{scrartcl}

\usepackage[utf8]{inputenc}
\usepackage[T1]{fontenc}
\usepackage{lmodern}
\usepackage{graphicx}
\usepackage[pdftex,hyperref,dvipsnames]{xcolor}
\usepackage{listings}
\usepackage[a4paper,lmargin={2cm},rmargin={2cm},tmargin={3.5cm},bmargin = {2.5cm},headheight = {4cm}]{geometry}
\usepackage{amsmath,amssymb,amstext,amsthm}
% \usepackage[lined,algonl,boxed]{algorithm2e} 
\usepackage{algpseudocode}
\usepackage{tikz}
\usepackage{pgfplots}
\usepackage{hyperref}
\usepackage{url}
\usepackage{tcolorbox}
\usepackage[inline]{enumitem}
\usepackage[headsepline]{scrlayer-scrpage} 

\usetikzlibrary{automata,positioning}

\usepgfplotslibrary{external}
%\tikzexternalize
\pgfplotsset{width=10cm,compat=1.9}
\pgfplotsset{every axis/.style={
    axis line style = thick,
    grid = both,
    tick pos = left,
    tick align = outside,
    tick style = {thick,black},
    fill = blue,
	line width = 2pt,
}}

\lstset{
  basicstyle=\ttfamily\small,
  numbers=left,
  numberstyle=\tiny,
  stepnumber=1,
  numbersep=8pt,
  frame=single,
  breaklines=true,
  columns=fullflexible
}

\lstset{
  language=C,
  basicstyle=\ttfamily\small,
  keywordstyle=\color{blue},
  commentstyle=\color{gray},
  stringstyle=\color{green!60!black},
  numbers=left,
  numberstyle=\tiny,
  stepnumber=1,
  numbersep=8pt,
  frame=single,
  breaklines=true,
  columns=fullflexible
}

\renewcommand{\theenumi}{(\alph{enumi})}
\renewcommand{\labelenumi}{\text{\theenumi}}
\renewcommand{\labelenumii}{(\roman{enumii})} 

\ohead{Weiting Wang (3855168)\\
       Nico Strohbach (3546914)\\
       Jan Theil (3794698)}

\ihead{Reinforcement Learning}

\newcounter{exnum}
\newcommand{\exercise}[1]{\section*{Problem \theexnum\stepcounter{exnum}: #1}}

\begin{document}

\setcounter{exnum}{1}

\exercise{Formulating Problem}

\begin{enumerate}
  \item Game of chess: The set of \textbf{states} is the discrete set of all possible board settings. This includes the position of all figures and the turn order of the oponents.
        The set of \textbf{actions} (for a specific state $S_i$) is the set of possible, valid moves based on $S_i$.
        The \textbf{reward signal} could be discrete: $+1$ for winning, $0$ for drawing and $-1$ for losing.
  \item Pick and place robot: The set of \textbf{states} is defined by the possible arm, motor positions of the robot and the localication of the movable object in question.
        Both of these (sub)properties are within $\mathbb{R}^3$, so this set is continuous.
        The set of \textbf{actions} is the set of possible robot movements / actions given the current state.
        The \textbf{reward signal} could be continuous: $+1$ for placing the object correctly and $-1$ for failing / breaking something. This can be scaled according to metrics like efficiency, shortest paths, etc.
  \item Drone: The set of \textbf{states} is again continuous and consists of the drones position, engine values and maybe sensor values.
        The set of \textbf{actions} is described by the possible engine control settings of the drone within a state.
        The \textbf{reward signal} can be $+1$ for succesfully stabilizing and $-1$ for crashing. Again this can be scaled by metrics like maximum unbalanced time, fuel efficiency, etc.
  \item Rollercoaster controller: The model should control a rollercoaster in a theme park.
        The set of \textbf{states} is the current position and velocity of the rollercoaster, this is continuous.
        The set of \textbf{actions} is discrete and consists of the actions to start and stop the rollercoaster and the discrete acceleration between $0$ and $100\%$.
        A possible \textbf{reward signal} would be $+1$ for not crashing, $-1$ for an accident and scaled by fun factor, etc.
\end{enumerate}

\exercise{Value Functions}

\begin{enumerate}
  \item In the setting of the $k$-armed bandit pulling an arm returns an immediate reward
        but it does not influence future rewards since in the next move any arm can be pulled again.
        So the future rewards are stochastically independent: $A_t = a$ has not causal effect on $R_{t + 2}, R_{t + 3}, \ldots$. Hence,
        $q(s, a) = \mathbb{E}[R_{t+1} + \gamma R_{t+2} + \gamma^2 R_{t+3} + \ldots | S_t = s, A_t = a] = C + \mathbb{E}[R_{t+1}| S_t = s, A_t = a]$.
        So, $q(s, a)$ is just offset by a constant $C$, which can be ignored when choosing an arm since every value is offset by the same $C$.
  \item \begin{align*}
        v_\pi(s) &= \mathbb{E}_\pi[G_t | S_t = s] \\
                 &= \mathbb{E}_\pi\left[ \sum_{i = 0}^\infty \gamma^i R_{t + i + 1 | S_t = s} \right] \\
                 &= \mathbb{E}_\pi [R_{t = 1} + \gamma G_{t+1} | S_t = s] \\
                 &= \sum_a \pi(a | s) \sum_{s'} \sum_r p(s', r | s, a)\left[ r + \gamma \mathbb{E}_\pi[G_{t + 1} | S_{t + 1} = s'] \right]
        \end{align*}
  \item $r(s, a, s')$ and $p(s' | s, a)$ are defined as:\\
        $r(s, a, s') = \sum_r r \cdot \dfrac{p(s', r | s, a)}{p(s' | s, a)} \implies \sum_r r \cdot p(s', r | s, a) = r(s, a, s') \cdot p(s' | s, a)$\\
        $p(s' | s, a) = \sum_r p(s', r | s, a)$\\
        We can substitute terms in $v_\pi(s)$ with these definitions:\\
        \begin{align*}
            v_\pi(s) &= \sum_{a} \pi(a | s) \sum_{s'} \sum_r p(s', r | s, a) \left[ r + \gamma v_\pi(s') \right] \\
                     &\text{factor out terms} \\
                     &= \sum_{a} \pi(a | s) \sum_{s'} \left[ \sum_r p(s', r | s, a) \cdot r + \sum_r p(s', r | s, a) \gamma v_\pi(s') \right] \\
                     &\text{use definitions from above} \\
                     &= \sum_{a} \pi(a | s) \sum_{s'} \left[ r(s, a, s') \cdot p(s' | s, a) + p(s' | s, a) \gamma v_\pi(s') \right] \\
                     &= \sum_{a} \pi(a | s) \sum_{s'} \left[ p(s' | s, a) \cdot (r(s, a, s') + \gamma v_\pi(s')) \right]
        \end{align*}

\end{enumerate}

\exercise{Bruteforce the Policy Space}

\begin{enumerate}
      \item The policy can choose between $|A|$ actions in each state. Therefore there are $|A|^{|S|}$ different policies.
      \item 
\end{enumerate}

\end{document}